\documentclass{article}
\usepackage{amsmath, amssymb, amsthm}
\usepackage{hyperref}
\usepackage{geometry}
\geometry{margin=1in}

\title{Erd\H{o}s Problem \#952}
\date{Last updated: 2025-08-31}
\author{Source: erdosproblems.com}

\begin{document}
\maketitle

\noindent\textbf{Status:} OPEN \quad \textbf{No prize}

\noindent\textbf{Tags:} number theory

\medskip
\noindent\textbf{Note:} Gaussian moat problem

\medskip

\noindent\textbf{Problem Statement:}

Is there an infinite sequence of distinct Gaussian primes $x_1,x_2,\ldots$ such that\[\lvert x_{n+1}-x_n\rvert \ll 1?\]




The Gaussian moat problem . This is not actually a problem of Erd\H{o}s, but has been erroneously attributed to him in the past. In \cite{Er77c} Erd\H{o}s recalls: 'The conjecture was told me by Motzkin at the Pasadena number theory meeting 1963 November and it was apparently raised by Basil Gordon and Motzkin. I naturally liked it very much and told it right away to many people, naturally attributing it to Motzkin, but this was later forgotten. Thus the problem is returned to its rightful owners.' In \cite{Er80} Erd\H{o}s writes 'the answer is almost certainly negative'.




References

[Er77c] Erd\H{o}s, Paul, Problems and results on combinatorial number theory. III . Number theory day (Proc. Conf., Rockefeller Univ., New York, 1976) (1977), 43-72.

[Er80] Erd\H{o}s, Paul, A survey of problems in combinatorial number theory . Ann. Discrete Math. (1980), 89-115.


\bigskip
\noindent\small{Source: \url{https://www.erdosproblems.com/952}}
\end{document}
